Suppose \(Y\) is flat over \(S\); then one has isomorphisms:
\[
J\mathcal{O}_{E\times_{S}Y}
  \isomfrom
  (J\mathcal{O}_{E})\otimes_{\mathcal{O}_{S}}\mathcal{O}_{Y}
  \isomto
  (J\mathcal{O}_{E})\otimes_{\mathcal{O}_{S_{0}}}\mathcal{O}_{Y_{0}}.
\]
Suppose moreover that \(E\) is of finite presentation over \(S\); then
\(\Omega^{1}_{E/S}\) is an \(\mathcal{O}_{E}\)-module of finite presentation
(\cf EGA~\(\mathrm{IV}_{4}\), 16.4.22) and hence, by EGA~\(0_{\mathrm{I}}\),
6.7.6, one has:
\[
\mathcal{H}\mathit{om}_{\mathcal{O}_{E\times_{S}Y}}
  \bigl(
    \Omega^{1}_{E/S}\otimes_{\mathcal{O}_{S}}\mathcal{O}_{Y},\,
    (J\mathcal{O}_{E})\otimes_{\mathcal{O}_{S}}\mathcal{O}_{Y}
  \bigr)
\simeq
\mathcal{H}\mathit{om}_{\mathcal{O}_{E}}
  \bigl(\Omega^{1}_{E/S},\, J\mathcal{O}_{E}\bigr)
  \otimes_{\mathcal{O}_{S}}\mathcal{O}_{Y}.
\]
Let us denote by \(\pi\colon E\to S\) and \(g\colon Y\to S\) the structural
morphisms; applying 0.13 to the diagram
\[
\xymatrix@C=4em@R=2.4em{
E \ar[d]_{\pi}
  & E\times_{S}Y \ar[l]_{g'} \ar[d]_{\pi'}
  \\
S
  & Y \ar[l]_{g}
\,,}
\]
and to the \(\mathcal{O}_{E}\)-module
\(\mathcal{F}
  =\mathcal{H}\mathit{om}_{\mathcal{O}_{E}}(\Omega^{1}_{E/S}, J\mathcal{O}_{E})\),
one obtains
\[
\Gamma(E\times_{S}Y,\, g'^{*}\mathcal{F})
  =\Gamma(Y,\, \pi'_{*}g'^{*}\mathcal{F})
  =\Gamma(Y,\, g^{*}\pi_{*}\mathcal{F})
  =W(\pi_{*}\mathcal{F})(Y).
\]
One has therefore shown that, if \(E\) is of finite presentation over \(S\),
one has
\begin{equation}
L'_{X}
  =W\bigl(
    \pi_{*}\mathcal{H}\mathit{om}_{\mathcal{O}_{E}}
      (\Omega^{1}_{E/S}, J\mathcal{O}_{E})
  \bigr)
\tag{0.14.1}
\end{equation}
on the category of \(S\)-schemes \emph{flat} over \(S\). Let us note moreover
that the module of which one takes the \(W\) is \emph{quasi-coherent}, by
EGA~I, 9.1.1 and 9.2.1.

{\renewcommand{\thefootnote}{(31)}\nde{One has added what follows.}}%
Let us denote by \(L_{0}\) the \(S_{0}\)-functor
\[
W\bigl(
  \pi_{0*}\mathcal{H}\mathit{om}_{\mathcal{O}_{E_{0}}}
    (\Omega^{1}_{E_{0}/S_{0}}, J\mathcal{O}_{E_{0}})
\bigr).
\]
Then, returning to the definition of \(L'_{X}(Y)\) and taking into account
the isomorphism
\[
J\mathcal{O}_{E\times_{S}Y}
  \simeq
  (J\mathcal{O}_{E})\otimes_{\mathcal{O}_{S_{0}}}\mathcal{O}_{Y_{0}},
\]
one obtains, by reasoning as above, that
\[
L'_{X}(Y)
  =L_{0}(Y_{0})
  =L_{0}(Y\times_{S}S_{0})
  =\Bigl(\prod_{S_{0}/S}L_{0}\Bigr)(Y).
\]
Thus, on the category of \(S\)-schemes \emph{flat} over \(S\), one has:
\[
L'_{X}
  =\prod_{S_{0}/S}
  W\bigl(
    \pi_{0*}\mathcal{H}\mathit{om}_{\mathcal{O}_{E_{0}}}
      (\Omega^{1}_{E_{0}/S_{0}}, J\mathcal{O}_{E_{0}})
  \bigr).
\]
It is not evident that the action of \(X\) on \(L'_{X}\) defined in 0.12
comes from an action of \(X_{0}=\Aut_{S_{0}}(E_{0})\) on \(L_{0}\); this is
however the case when, moreover, \(E\) is \emph{flat} over \(S\).

Indeed, one has in this case canonical isomorphisms:
\begin{align*}
J\mathcal{O}_{E}
  &\simeq J\otimes_{\mathcal{O}_{S}}\mathcal{O}_{E}
  \simeq J\otimes_{\mathcal{O}_{S_{0}}}\mathcal{O}_{E_{0}}.\\
L_{0}
  &\simeq W\bigl(
    \pi_{0*}\mathcal{H}\mathit{om}_{\mathcal{O}_{E_{0}}}
      (\Omega^{1}_{E_{0}/S_{0}},
        J\otimes_{\mathcal{O}_{S_{0}}}\mathcal{O}_{E_{0}})
  \bigr),
\end{align*}
\begin{equation}
L'_{X}
  =\prod_{S_{0}/S}
  W\bigl(
    \pi_{0*}\mathcal{H}\mathit{om}_{\mathcal{O}_{E_{0}}}
      (\Omega^{1}_{E_{0}/S_{0}},
        J\otimes_{\mathcal{O}_{S_{0}}}\mathcal{O}_{E_{0}})
  \bigr).
\tag{0.14.2}
\end{equation}
Then, for every \(S_{0}\)-scheme \(T\), one has
\[
L_{0}(T)
  \simeq
  \Hom_{\mathcal{O}_{E_{0}\times_{S_{0}}T}}
    \bigl(
      \Omega^{1}_{E_{0}\times_{S_{0}}T/T},\,
      J\otimes_{\mathcal{O}_{S_{0}}}\mathcal{O}_{E_{0}\times_{S_{0}}T}
    \bigr)
\]
and \(\Hom_{S_{0}}(T, X_{0})=\Aut_{T}(E_{0}\times_{S_{0}}T)\) acts by
transport of structure on \(L_{0}(T)\), in a manner functorial in \(T\);
finally, for every \(S\)-scheme \(Y\) flat over \(S\), the action by transport
of structure of \(\Hom_{S}(Y, X)=\Aut_{Y}(E\times_{S}Y)\) on
\(L'_{X}(Y)=L_{0}(Y_{0})\) factors through
\(\Aut_{Y_{0}}(E_{0}\times_{S_{0}}Y_{0})\).

Let us finally extract from SGA~1 III the following two propositions.

\begin{proposition}{0.15}\label{III.0.15}
(SGA~1 III, 6.8)%
{\renewcommand{\thefootnote}{(32)}\nde{This \og global\fg\ lifting result
uses the \og local\fg\ lifting result of \loccit, 4.1, which is stated for
\(S\) locally noetherian; see EGA~\(\mathrm{IV}_{4}\) 18.1.1 for the general
case.}}
For every \(S_{J}\)-scheme \(Y\) smooth over \(S_{J}\) and affine, there
exists an \(S\)-scheme \(X\) smooth over \(S\) such that
\(X\times_{S}S_{J}\simeq Y\), and such an \(X\) is unique up to
(non-unique) isomorphism.
\end{proposition}

\begin{proposition}{0.16}\label{III.0.16}
(SGA~1 III, 5.5)%
{\renewcommand{\thefootnote}{(33)}\nde{This is an immediate consequence of
the definition of (formally) \og smooth\fg\ adopted in
EGA~\(\mathrm{IV}_{4}\), 17.3.1 (and 17.1.1).}}
Let \(X\) be an \(S\)-scheme smooth over \(S\). For every affine \(S\)-scheme
\(Y\), the canonical map
\[
p_{X}(Y)\colon
  \Hom_{S}(Y, X)
  \longrightarrow
  \Hom_{S}(Y, X^{+})
  =\Hom_{S_{J}}(Y_{J}, X_{J})
\]
is surjective.
\end{proposition}

\begin{corollary}{0.17}\label{III.0.17}
Let \(E\) be an \(S\)-scheme smooth over \(S\) and affine; let us write
\(X=\Aut_{S}(E)\). For every affine \(S\)-scheme \(Y\), the canonical map
\[
\Aut_{Y}(E\times_{S}Y)
  =\Hom_{S}(Y, X)
  \longrightarrow
  \Hom_{S}(Y, X^{+})
  =\Aut_{Y_{J}}(E_{J}\times_{S_{J}}Y_{J})
\]
is surjective.
\end{corollary}

\begin{proof}
Indeed, \(Y\times_{S}E\) is affine over \(Y\), itself affine, hence affine.
Applying 0.16, one deduces from this that every \(S_{J}\)-morphism
\(Y_{J}\times_{S_{J}}E_{J}\to E_{J}\) extends to an \(S\)-morphism
\(Y\times_{S}E\to E\).
{\renewcommand{\thefootnote}{(34)}\nde{One has slightly modified the
original, so as to be exactly under the hypotheses of 0.7.}}
In other terms, every \(Y_{J}\)-endomorphism of
\(Y_{J}\times_{S_{J}}E_{J}\) lifts to a \(Y\)-endomorphism of
\(Y\times_{S}E\). Then, 0.7~a) shows that every
\emph{\(Y_{J}\)-automorphism} of \(Y_{J}\times_{S_{J}}E_{J}\) lifts to a
\emph{\(Y\)-automorphism} of \(Y\times_{S}E\), which is the announced
property.
\end{proof}

\sgasection{1}{Extensions and cohomology}
\label{III.sec.1}

\numpara{1.1}
Let \(\mathcal{C}\) be a category stable under fiber
products.%
{\renewcommand{\thefootnote}{(35)}\nde{One has added the hypothesis that
\(\mathcal{C}\) be stable under fiber products (which is the case in the
applications where \(\mathcal{C}\) is the category of schemes \Sch\ or that
of the functors \(\Sch^{\circ}\to\Ens\)). If one omits this hypothesis, one
must suppose in the sequel that the fiber products
\(G\times_{S}\cdots\times_{S}G\) are representable.}}
Let \(S\) be an object of \(\mathcal{C}\), \(G\) an \(S\)-group
(\emph{representable}) and \(F\) a commutative group \(S\)-functor on which
\(G\) operates. One has defined in~I, 5.1 the cohomology groups
\(H^{n}(G, F)\). One recalls that these are the homology groups of a complex
denoted \(C^{*}(G, F)\) where, writing
\((G/S)^{n}=G\times_{S}\cdots\times_{S}G\) (\(n\) factors),
\[
C^{n}(G, F)=\Hom_{S}((G/S)^{n}, F).
\]
As \(G\) and hence\textsuperscript{(35)} the \((G/S)^{n}\) are representable,
one also has
\[
C^{n}(G, F)=F((G/S)^{n})\,;
\]
from this, and from the definition of the boundary operator, one sees that
the complex \(C^{*}(G, F)\) depends only on the restriction of \(F\) to the
full subcategory of \(\mathcal{C}_{/S}\) whose objects are the cartesian
powers of \(G\) over \(S\). Consequently, one has the

\begin{lemma}{1.1.1}\label{III.1.1.1}
Let \(\mathcal{C}\) be a category stable under fiber
products\textsuperscript{(35)}, \(S\) an object of \(\mathcal{C}\), \(G\) a
representable \(S\)-group. Let us denote by \(\mathcal{C}(G)\) the full
subcategory of \(\mathcal{C}_{/S}\) whose objects are the cartesian powers
of \(G\) over \(S\). Let \(F\) and \(F'\) be two commutative group
\(S\)-functors on which \(G\) operates. If \(F\) and \(F'\) have the same
restriction to \(\mathcal{C}(G)\), one has a canonical isomorphism
\[
H^{*}(G, F)\isomto H^{*}(G, F').
\]
\end{lemma}

\par\addvspace{0.55\baselineskip}%
\noindent\textbf{1.1.2. Cohomology and restriction of scalars.}\ ---%
\label{III.ssec.1.1.2}
{\renewcommand{\thefootnote}{(36)}\nde{One has added the title of this
paragraph.}}
Let us state another comparison result. Let now \(T\to S\) be a morphism of
\(\mathcal{C}\). If \(F\) is a commutative group \(T\)-functor, then the
functor obtained by \og restriction of scalars\fg\ (\cf Exp.~II, 1)
\[
F_{1}=\prod_{T/S}F
\]
is a commutative group \(S\)-functor and one has a morphism of group
\(S\)-functors
\[
u\colon
  \prod_{T/S}\Aut_{\mathrm{T}\text{-gr.}}(F)
  \longrightarrow
  \Aut_{\mathrm{S}\text{-gr.}}(F_{1}).
\]%
{\renewcommand{\thefootnote}{(37)}\nde{Indeed, let \(S'\to S\) and
\(\alpha\in\Aut_{(T\times_{S}S')\text{-gr.}}(F_{T\times_{S}S'})\); for every
\(U\to T\times_{S}S'\), let us denote by
\(\alpha_{U}\in\Aut_{\mathrm{gr.}}(F_{T\times_{S}S'}(U))\) the element
defined by \(\alpha\); then \(u(S')(\alpha)\) is the element \(\beta\) of
\(\Aut_{\mathrm{S}\text{-gr.}}(F_{1})(S')
  =\Aut_{\mathrm{S}'\text{-gr.}}
    \bigl(\prod_{T\times_{S}S'/S'}F_{T\times_{S}S'}\bigr)\)
such that \(\beta_{S''}=\alpha_{T\times_{S}S''}\) for every
\(S''\to S'\).}}
Let now \(G\) be a group \(S\)-functor and let
\[
G_{T}\longrightarrow\Aut_{\mathrm{T}\text{-gr.}}(F)
\]
be an operation of \(G_{T}\) on \(F\). By definition of the functor
\(\prod_{T/S}\), one deduces from this a morphism of group \(S\)-functors
\[
G\longrightarrow\prod_{T/S}\Aut_{\mathrm{T}\text{-gr.}}(F),
\]
whence, by composition with \(u\), an operation of \(G\) on
\(F_{1}=\prod_{T/S}F\).%
{\renewcommand{\thefootnote}{(38)}\nde{Explicitly, for every \(S'\to S\),
the operation of \(G(S')\) on \(F_{1}(S')=F(T\times_{S}S')\) is given by the
operation of \(G_{T}(T\times_{S}S')=G(T\times_{S}S')\) on
\(F(T\times_{S}S')\) and the morphism of groups
\(G(S')\to G(T\times_{S}S')\) corresponding to the projection
\(T\times_{S}S'\to S'\). One may also say that the operation of \(G_{T}\) on
\(F\) is given by a morphism \(G_{T}\times_{T}F\to F\); applying the functor
\(\prod_{T/S}\) and denoting by \(G_{1}\) the \(S\)-group
\(\prod_{T/S}G_{T}\), one obtains, taking into account II~1.2, a morphism
\(G_{1}\times_{S}F_{1}\to F_{1}\) which defines an operation of \(G_{1}\) on
\(F\). The operation of \(G\) is then obtained via the natural morphism
\(G\to G_{1}\), which corresponds by adjunction to \(\mathrm{id}_{G}\).}}

\begin{lemma}{1.1.2}\label{III.1.1.2}
Under the preceding conditions, one has a canonical isomorphism
\[
H^{*}\Bigl(G, \prod_{T/S}F\Bigr)\simeq H^{*}(G_{T}, F).
\]
\end{lemma}

\begin{proof}
Indeed, by the definition of the cohomology, the standard complexes are
canonically isomorphic.
\end{proof}

\par\addvspace{0.55\baselineskip}%
\noindent\textbf{1.2. Lifting of morphisms of groups.}\ ---%
\label{III.ssec.1.2}
{\renewcommand{\thefootnote}{(39)}\nde{One has added this title.}}
Following the general principles, one sets down the following definition:

\begin{definition}{1.2.1}\label{III.1.2.1}
Let \(1\longrightarrow M\xrightarrow{u}E\xrightarrow{v}G\) be a sequence of
morphisms of \(\widehat{\mathcal{C}}\)-groups. One says that it is
\emph{exact} if the following equivalent conditions are verified:
\begin{enumeratei}
\item for every \(S\in\Ob\mathcal{C}\), the sequence of ordinary groups
below is exact:
\[
1\longrightarrow M(S)\xrightarrow{u(S)}E(S)\xrightarrow{v(S)}G(S)
\]
\item for every object \(H\) of \(\widehat{\mathcal{C}}\), the sequence of
ordinary groups below is exact:
\[
1\longrightarrow
  \Hom(H, M)\xrightarrow{u(H)}\Hom(H, E)\xrightarrow{v(H)}\Hom(H, G)
\]
\end{enumeratei}
\end{definition}

Taking in particular \(H=G\) in~(ii), one sees that the set of sections of
\(v\) (not respecting \emph{a priori} the group structures) is empty or
principal homogeneous under \(\Hom(G, M)\). Let us suppose it non-empty; let
therefore
\[
s\colon G\longrightarrow E
\]
be a section of \(v\). Then for every \(S\in\Ob\mathcal{C}\) and every
\(x\in G(S)\), the element \(s(x)\) of \(E(S)\) defines an inner automorphism
of \(E_{S}\) which normalizes \(M_{S}\) (more correctly the image of
\(M_{S}\) under \(u_{S}\)), hence an automorphism of \(M_{S}\).

\begin{scholium}{1.2.1.1}\label{III.1.2.1.1}
{\renewcommand{\thefootnote}{(40)}\nde{One has added the numbering 1.2.1.1
and 1.2.1.2 in order to bring out the notions which are introduced
there.}}%
If \(M\) is commutative, one sees \og set-theoretically\fg\ that this
automorphism does not depend on the chosen section, but only on \(x\), and
that it depends multiplicatively on it. In summary, to every exact sequence
\[
(E)\qquad
1\longrightarrow M\xrightarrow{u}E\xrightarrow{v}G
\]
such that \(M\) is commutative and that \(v\) possesses a section, there is
associated a morphism of \(\widehat{\mathcal{C}}\)-groups
\[
G\longrightarrow\Aut_{\widehat{\mathcal{C}}\text{-gr.}}(M)
\]
which one calls the operation of \(G\) on \(M\) defined by the extension
\((E)\).
\end{scholium}

\begin{definition}{1.2.1.2}\label{III.1.2.1.2}
One has seen in~I, 2.3.7 that \(v\) possesses a section which is a morphism
of \(\widehat{\mathcal{C}}\)-groups if and only if the extension \((E)\) is
isomorphic (\og as an extension\fg) to the semidirect product of \(M\) by
\(G\) relative to the preceding operation. Such a section of \(v\) will be
called a \emph{section of the extension \((E)\)}, or simply a
\emph{section of \((E)\)}.

If \(s\) is a section of \((E)\) and if
\(m\in\Gamma(M)\simeq\Ker(\Gamma(E)\to\Gamma(G))\) (for the definition of
\(\Gamma\), see~I, 1.2), then the morphism \(G\to E\) defined
by%
{\renewcommand{\thefootnote}{(41)}\nde{In the right-hand term, one has
corrected \(x\) to \(s(x)\).}}
\[
x\mapsto u(m)\,s(x)\,u(m)^{-1}
\]
is likewise a section of \((E)\), said to be \emph{deduced from \(s\) by the
inner automorphism defined by \(m\) (or by \(u(m)\))}.
\end{definition}

\begin{lemma}{1.2.2}\label{III.1.2.2}
Let \((E)\colon 1\longrightarrow M\xrightarrow{u}E\xrightarrow{v}G\) be an
exact sequence of \(\widehat{\mathcal{C}}\)-groups such that \(M\) is
commutative and that \(v\) possesses a section. Let us make \(G\) operate on
\(M\) in the manner defined by \((E)\).
\begin{enumeratei}
\item The extension \((E)\) canonically defines a class
\(c(E)\in H^{2}(G, M)\) whose vanishing is necessary and sufficient for the
existence of a section of \((E)\).
\item If \(c(E)=0\), the set of sections of \((E)\) is principal homogeneous
under the group \(Z^{1}(G, M)\), and the set of sections of \((E)\) modulo
the action of the inner automorphisms defined by the elements of
\(\Gamma(M)\) is principal homogeneous under the group \(H^{1}(G, M)\).
\item {\renewcommand{\thefootnote}{(42)}\nde{One has added this point, the
counterpart of 1.2.4~(iii). Let us note on the other hand that the
assertion is valid for \emph{every section} of \(v\), \cf the proof.}}%
Let \(s\) be a section of \((E)\); the set of the conjugates of \(s\) by the
inner automorphisms defined by \(\Gamma(M)\) is in bijection with
\(\Gamma(M)/H^{0}(G, M)\).
\end{enumeratei}
\end{lemma}

\begin{proof}
The proof is carried out exactly as in the case of ordinary groups, the
fact that one starts from a section of \(v\) ensuring the functoriality of
the set-theoretic computations. Let us briefly indicate the principal steps
of the proof.

\textbf{a)} To every section \(s\) of \(v\) one associates the morphism
\[
Ds\colon G\times G\longrightarrow M,
\]
defined set-theoretically by
\[
u\bigl(Ds(x, y)\bigr)=s(xy)\,s(y)^{-1}\,s(x)^{-1}.
\]
One shows that \(Ds\) is a \(2\)-cocycle by the following
computation.%
{\renewcommand{\thefootnote}{(43)}\nde{One has corrected the original, so
as to make it compatible with I, 5.1.}}
By the definition of the differential of the standard complex (I, 5.1), one
has:
\[
(\partial^{2}Ds)(x, y, z)
  =\bigl(s(x)\,Ds(y, z)\,s(x)^{-1}\bigr)
  \cdot Ds(x, y)^{-1}
  \cdot Ds(xy, z)^{-1}
  \cdot Ds(x, yz);
\]
it suffices to substitute the definition of \(Ds\) into this formula in
order to find (without using any commutativity) \(Ds\in Z^{2}(G, M)\).

\textbf{b)} If \(s\) and \(s'\) are two sections of \(v\), there exists
\(f\colon G\to M\) such that \(s(x)=f(x)s'(x)\). One then has
\begin{align*}
Ds'(x, y)
  &=f^{-1}(xy)\,Ds(x, y)\,s(x)\,f(y)\,s(x)^{-1}\,f(x),\\
  &=Ds(x, y)\cdot f^{-1}(xy)
    \cdot\bigl(s(x)\,f(y)\,s(x)^{-1}\bigr)\cdot f(x),
\end{align*}
that is,
\[
Ds'=Ds\cdot\partial^{1}f.
\]
{\renewcommand{\thefootnote}{(44)}\nde{One has added the sentence which
follows.}}
This shows that the class of \(Ds\) in \(H^{2}(G, M)\) does not depend on the
chosen section \(s\) of \(v\); it is the class \(c(E)\) of the extension
\((E)\).

\textbf{c)} Let \(s\) and \(s'\) be two sections of \(v\) and let
\(m\in\Gamma(M)\). Then, the equality \(s(x)=m^{-1}s'(x)m\) (for every
\(x\in G(S)\), \(S\in\Ob\mathcal{C}\)) is equivalent to
\[
s(x)=m^{-1}s'(x)\,m\,s'(x)^{-1}s'(x),
\qquad\ie\qquad
s=\partial^{0}m\cdot s'.
\]
In particular, the stabilizer of \(s\) in \(\Gamma(M)\) is the subgroup of
those \(m\in\Gamma(M)\) such that \(\partial^{0}m=e_{M}\), \ie the subgroup
\(H^{0}(G, M)\). This already proves~(iii).

\textbf{d)} The reasoning is now the usual one:%
{\renewcommand{\thefootnote}{(45)}\nde{One has set out in detail what
follows.}}
Let \(s_{0}\) be an arbitrary section of \(v\); there exists a section
\(s\), necessarily of the form \(s=f\cdot s_{0}\), which is a morphism of
groups, \ie which satisfies \(Ds=0\), if and only if
\((Ds_{0})^{-1}=\partial^{1}f\), \ie, if and only if the class \(c(E)\) is
zero. This proves~(i).

In this case, the set of sections of \((E)\) is formed of the sections
\(s'=h\cdot s\), where \(h\colon G\to M\) satisfies \(\partial^{1}h=0\),
\ie \(h\in Z^{1}(G, M)\). Moreover, by point~c), two such sections
\(h_{1}\cdot s\) and \(h_{2}\cdot s\) are conjugate under \(\Gamma(M)\) if
and only if \(h_{1}\) and \(h_{2}\) have the same image in
\(H^{1}(G, M)\). This proves~(ii).
\end{proof}

Let still
\[
(E)\qquad
1\longrightarrow M\xrightarrow{u}E\xrightarrow{v}G
\]
be an exact sequence of \(\widehat{\mathcal{C}}\)-groups with \(M\)
\emph{commutative}. Let
\[
f\colon H\longrightarrow G
\]
be a morphism of \(\widehat{\mathcal{C}}\)-groups.
